%=======================================================================
% cap6_conclusiones.tex — Conclusiones, desafíos y trabajo futuro
%=======================================================================

%-----------------------------------------------------------------------
\section{Respuesta a las preguntas de investigación}\label{sec:respuestas}

\textbf{PI1. ¿Puede un sistema multimodal identificar aves chilenas mejor
que los clasificadores unimodales (solo audio o solo imagen) y que un modelo de IA generativa multimodal?} Sí, en ambas comparaciones. La \textit{late fusion} de audio e imagen alcanza 0,945 $\pm$ 0,004 de
\textit{accuracy} sobre 53 especies (intervalo de confianza al
\SI{95}{\percent} de 0,924 a 0,959), 8,0 puntos más que el mejor
clasificador unimodal sobre los mismos pares y casi 10 sobre la
meta comprometida, y lo hace con datos no emparejados de ciencia ciudadana,
un escenario que la literatura previa no había explorado. Cada modalidad
por separado también cumple o roza la meta: audio 0,858 (CNN) y 0,870 (AST)
sobre 55 especies, imagen 0,848 sobre 53. Frente al modelo de IA generativa multimodal, el sistema especializado gana con claridad: Claude Opus 5.5 alcanza 0,803 en \textit{zero-shot} y 0,817 en
\textit{few-shot} sobre las mismas fotografías, 3 a 4,5 puntos menos que la
rama de imagen y 13 menos que la fusión, y falla precisamente en las
especies endémicas chilenas (F1 0,498 y 0,567 frente a 0,940 de la
fusión).

\textbf{PI2. ¿Puede un SDM basado en aprendizaje automático superar a
MaxEnt?} Sí, pero no con redes neuronales. Random Forest y gradient
boosting superan a MaxEnt en AUC-ROC (0,884 y 0,883 frente a 0,863) en 46 y 44
de las 54 especies, mientras que la red neuronal conjunta queda por debajo
(0,856) y la incorporación de parches de cobertura de suelo y
\textit{embeddings} de imágenes satelitales de un modelo fundacional degrada
a todos los modelos. La hipótesis ``aprendizaje
profundo $>$ MaxEnt'' se refuta con evidencia. La hipótesis ``aprendizaje
automático $>$ MaxEnt'' se confirma con un método más simple, pero con una
condición importante: bajo validación espacial por bloques, que mide la
extrapolación a zonas sin registros, todos los modelos caen a 0,80--0,82 y
ninguno supera a MaxEnt. El aprendizaje automático gana al interpolar, no al
extrapolar.

\textbf{PI3. ¿Cómo afecta la escasez de datos de las especies amenazadas y
qué estrategias la mitigan?} De forma distinta en cada módulo. En imagen,
donde los plumajes de congéneres se parecen, las especies con menos
fotografías son las que fallan (F1 medio 0,756 en amenazadas frente a 0,842
en comunes) y dos quedaron fuera del clasificador. En audio, la escasez se
mitigó con grabaciones de otras poblaciones, aumentación en el espectrograma
y pesos por clase, y las amenazadas rinden incluso mejor que las comunes
(0,906 frente a 0,838) porque sus vocalizaciones son distintivas. El precio
es la dependencia de audio extranjero y conjuntos de prueba de tres a ocho
grabaciones. En distribución, los pocos registros de las especies de rango
restringido bastan para modelarlas casi perfectamente (AUC 0,967). La fusión
multimodal es, en sí misma, la mitigación más efectiva: cubre las especies
que una modalidad no puede clasificar y corrige los errores de la otra.

%-----------------------------------------------------------------------
\section{Discusión}\label{sec:discusion}

\textbf{Por qué funciona la fusión.} La complementariedad no es casual. El
audio separa especies por estructura temporal y espectral del canto, que en
aves es un carácter de reconocimiento de especie. La imagen separa por forma
y plumaje. Los pares que confunde una modalidad (\textit{Cinclodes} y
\textit{Phrygilus} en audio, \textit{Scytalopus} en imagen) son distintos de
los que confunde la otra, y la \textit{late fusion} solo necesita que una de las dos
esté segura. Que la \textit{late fusion} y la \textit{early fusion} empaten y
la \textit{intermediate fusion} pierda indica que, con \textit{backbones}
fuertes y pocos pares, lo simple gana: la forma más barata de combinar
evidencia es también la mejor, y la \textit{late fusion} tiene la ventaja de
operar con una sola modalidad.

\textbf{Por qué los árboles ganan al aprendizaje profundo en el SDM.} El
problema de distribución, tal como se formuló, es tabular: seis covariables
y unos cientos de presencias por especie. En ese régimen los ensambles de
árboles tienen sesgos inductivos adecuados (invariancia a la escala de las
variables, fronteras de decisión irregulares, robustez a variables
irrelevantes) y las redes necesitan más datos o regularización
\citep{grinsztajn2022trees}. La red conjunta comparte información entre
especies, pero esa ventaja no compensa la pérdida frente a 300 árboles por
especie. Los experimentos satelitales muestran además que la señal
ambiental relevante ya está en las covariables clásicas: el NDVI, el clima y
la elevación resumen lo que la cobertura de suelo y el \textit{embedding}
satelital repiten con más ruido. Y la validación espacial agrega
un matiz: la ventaja de los árboles es de interpolación. Al extrapolar a
zonas sin registros, la respuesta suave y regularizada de MaxEnt generaliza
igual de bien. Este resultado es informativo para la práctica: con buenas
covariables, un Random Forest es la elección razonable para describir la
distribución donde hay datos, MaxEnt sigue siendo tan bueno como cualquiera
para proyectarla a zonas sin muestreo, y la complejidad adicional debe
justificarse con una ganancia medida, no supuesta.

\textbf{El valor de la agregación y de los protocolos operacionales.} La
métrica por ventana única subestima el desempeño real en cuatro puntos.
Evaluar como se usará el sistema (la grabación completa, la fotografía con
TTA, el par audio-imagen) no solo da números más altos sino más honestos
respecto del caso de uso. Lo mismo vale para la comparación de costo: el AST
mejora un punto a la CNN con veinte veces el cómputo, una diferencia que
decide qué modelo conviene desplegar en un dispositivo de campo.

\textbf{Relación con la literatura.} La ganancia de la fusión coincide con lo
observado en SSW60, donde combinar audio e imagen supera a cada modalidad
\citep{vanhorn2022ssw60}, y la extiende a datos no emparejados de ciencia
ciudadana. En distribución, el resultado es coherente con el
\textit{benchmark} de \citet{valavi2022benchmark}, donde MaxEnt y los
ensambles de árboles quedan entre los mejores métodos con diferencias
pequeñas, y con la advertencia de \citet{roberts2017crossvalidation}: la
partición aleatoria sobreestima el desempeño cuando los datos tienen
estructura espacial, aquí entre cinco y ocho puntos de AUC. Por último, el
desempeño de la rama de audio sobre especies chilenas, frente a la cobertura
moderada de BirdNET en Sudamérica \citep{funosas2026birdnet}, respalda la
estrategia de entrenar clasificadores regionales con datos locales, aunque
las dos evaluaciones no son directamente comparables.

%-----------------------------------------------------------------------
\section{Desafíos y limitaciones}\label{sec:limitaciones}

\begin{enumerate}[leftmargin=*, itemsep=3pt]
  \item \textbf{Evaluación de la fusión con pares sintéticos.} Los pares
        audio-imagen se construyen entre individuos distintos de la misma
        especie. El protocolo mide la ganancia de combinar evidencia, que es
        el caso de uso previsto, pero no captura correlaciones entre el canto
        y la apariencia de un mismo individuo ni condiciones de campo donde
        ambas señales se degradan a la vez (un ave lejana, grabada y
        fotografiada con ruido).
  \item \textbf{Conjuntos de prueba pequeños para las especies raras.} Con
        tres grabaciones o dos fotografías de prueba, el F1 de una especie
        amenazada puede oscilar entre 0 y 1 por un solo ejemplo. Las métricas
        globales son robustas (la repetición del pareo da una desviación
        estándar de 0,004 en la fusión), pero las de especies individuales
        con pocos datos deben leerse como indicativas: el F1 mínimo dentro del
        grupo de amenazadas baja a 0,48 por una sola especie con dos
        fotografías.
  \item \textbf{Origen de las grabaciones.} El \SI{36}{\percent} del audio
        proviene de otros países y cuatro especies no tienen audio chileno.
        El conjunto de prueba es mixto. El análisis por origen no muestra
        penalización (la CNN obtiene 0,847 en las grabaciones chilenas y 0,877
        en las extranjeras, mientras que el AST obtiene 0,871 y 0,870), pero la limitación
        persiste en otro sentido: cuatro especies solo pueden evaluarse con
        audio extranjero, y nada garantiza ese desempeño sobre paisajes
        sonoros grabados en terreno en Chile.
  \item \textbf{Sesgos de los registros de presencia.} Los avistamientos de
        GBIF carecen de esfuerzo de observación, se concentran en 2024--2026 y
        en el verano austral, y siguen la geografía de los observadores. El
        fondo de grupo objetivo corrige el sesgo espacial, pero la estacionalidad
        aprendida por el modelo mezcla fenología con esfuerzo, y el tope de
        \num{2000} registros por especie descarta historia útil.
  \item \textbf{Validación espacial.} La comparación principal del SDM usa una
        partición aleatoria, optimista ante la autocorrelación espacial. La
        validación por bloques cuantifica esa diferencia: entre cinco y ocho
        puntos de AUC, y borra la ventaja de los árboles. Por eso ambos
        protocolos se reportan juntos y la conclusión sobre el SDM se formula
        con esa salvedad.
  \item \textbf{Covariables estáticas.} WorldClim y el NDVI medio anual
        describen el clima y la vegetación promedio, no el año ni el mes de
        cada observación. El pareo anual del \textit{embedding} satelital
        mostró que ese desajuste pesa poco, pero una serie temporal de NDVI a
        mayor resolución sigue siendo la vía natural para un modelo
        verdaderamente espaciotemporal.
  \item \textbf{Anotación débil de Xeno-canto.} Cada grabación lleva una
        especie focal, pero puede contener otras de fondo y tramos sin
        vocalización. La agregación por grabación mitiga el efecto, pero no lo
        elimina, y las ventanas de entrenamiento siguen siendo ruidosas.
  \item \textbf{Cómputo local.} Las arquitecturas de imagen más fuertes
        (ConvNeXt y ViT-B con preentrenamiento en ImageNet-21k a
        \SI{384}{px}) quedaron fuera porque su tiempo de entrenamiento era
        inviable en el equipo; el ViT-S/16 que sí se entrenó, a \SI{224}{px},
        quedó entre seis y ocho puntos por debajo de EfficientNet-B3 según el
        protocolo, de modo que sigue abierta la pregunta de si la atención
        global ayuda con mayor resolución y más cómputo; con GPU dedicada es
        probable que la rama de imagen mejore. El AST, por su parte, se
        entrenó con lotes de 8 por la misma restricción.
  \item \textbf{Comparación con el modelo generativo.} El resultado depende del
        modelo y la versión consultados, del \textit{prompt} y de que las
        fotografías de iNaturalist puedan haber formado parte de su
        entrenamiento. Se reporta como referencia, no como \textit{benchmark}
        estable.
\end{enumerate}

\section{Aspectos éticos}\label{sec:etica}

El proyecto no involucra personas ni datos personales, por lo que no requirió
la aprobación de un comité de ética. Aun así, el uso de datos de ciencia
ciudadana y la publicación de modelos sobre especies amenazadas plantean
cuatro consideraciones.

\textbf{Origen y licencias de los datos.} Todas las fuentes son públicas y se
usaron según sus condiciones. Las grabaciones de Xeno-canto y las fotografías
de iNaturalist llevan licencias Creative Commons elegidas por sus autores y
se emplearon solo para entrenar y evaluar modelos, sin redistribuirlas. Los
registros de eBird e iNaturalist se obtuvieron de GBIF y se citan con el
identificador de cada conjunto \citep{ebird2026gbif, inat2026gbif}. Las capas
ambientales (WorldClim, MODIS, SRTM, WorldCover) son de acceso abierto. De
cada registro se usan solo la especie, las coordenadas y la fecha; no se
procesa la identidad de los observadores, y las plataformas ya ofuscan las
localidades de las especies que consideran sensibles antes de publicarlas.

\textbf{Riesgo de mal uso.} Un mapa de probabilidad de presencia de una
especie amenazada puede orientar su captura o su tráfico; los loros, como el
tricahue, son el caso más conocido en Chile. Tres decisiones acotan ese
riesgo. Los mapas se calculan sobre una grilla de \SI{0,1}{\degree} (unos
\SI{10}{\kilo\metre}) y describen la idoneidad del hábitat, no la ubicación
de individuos ni de nidos. El prototipo no muestra ni entrega los registros
individuales que alimentan los modelos. Y en un despliegue público los mapas
de las especies más sensibles deberían publicarse a menor resolución o
restringirse a instituciones de conservación, como hacen eBird e iNaturalist
con sus datos sensibles.

\textbf{Decisiones automatizadas.} El clasificador se equivoca más en las
especies con pocos datos y confunde congéneres de plumaje o canto parecidos
(Secciones~\ref{sec:res-audio} y \ref{sec:res-imagen}). Si sus predicciones
alimentaran una base de datos sin revisión, introducirían registros falsos
justamente en las especies que más importan. Por eso el sistema se plantea
como apoyo a la identificación experta: el prototipo muestra las especies
candidatas con su probabilidad en lugar de una etiqueta única, y cualquier
uso para registros oficiales debe pasar por validación humana.

\textbf{Sesgos de los datos.} Los registros ciudadanos siguen a los
observadores: se concentran en la zona central, en verano y en los últimos
años (\secref{sec:limitaciones}). Un modelo entrenado con ellos describe
mejor las zonas visitadas que las remotas, y las prioridades de monitoreo
que se deriven de sus mapas podrían reforzar ese sesgo. El fondo de grupo
objetivo y la validación por bloques lo acotan y lo cuantifican, pero no lo
eliminan, y así debe leerse cualquier mapa del sistema. Por último, las
fotografías enviadas al modelo generativo comercial son públicas y con
licencia abierta, y no se envió información sobre los observadores.

%-----------------------------------------------------------------------
\section{Lecciones aprendidas y recomendaciones}\label{sec:lecciones}

Más allá de los resultados, el proyecto deja lecciones que pueden servir a
otros estudiantes del programa.

\begin{itemize}[leftmargin=*, itemsep=3pt]
  \item \textbf{Pedir los datos restringidos el primer día.} La base
        histórica de eBird se solicitó en julio y no llegó a tiempo, lo que
        obligó a reemplazar las no detecciones reales por pseudoausencias.
        Un acceso que depende de terceros debe gestionarse al inicio y con
        un plan alternativo.
  \item \textbf{Fijar los \textit{baselines} y el protocolo de evaluación
        antes de entrenar.} Decidir de antemano contra qué se compara el
        sistema (MaxEnt, \textit{frozen embeddings}, un modelo generativo) y
        cómo se mide (ventana única o agregado, recorte central o TTA, pareo
        por especie) evitó ajustar la vara después de ver los resultados.
  \item \textbf{Evaluar como se usará el sistema.} La métrica por ventana
        única subestimaba el desempeño en cuatro puntos; la grabación
        completa, la fotografía con TTA y el par audio-imagen son los
        escenarios reales de uso.
  \item \textbf{Validar espacialmente desde el inicio.} La partición
        aleatoria daba a los árboles una ventaja de dos puntos sobre MaxEnt
        que la validación por bloques borró. Con datos georreferenciados, la
        validación espacial debe ser el protocolo por defecto y no una
        comprobación final.
  \item \textbf{No presuponer que el aprendizaje profundo gana.} En datos
        tabulares la red neuronal perdió frente a Random Forest y frente a
        MaxEnt. La complejidad de un modelo se justifica con una ganancia
        medida.
  \item \textbf{Medir el costo junto con la métrica.} El AST gana un punto a
        la CNN con veinte veces el cómputo; sin ese dato la decisión de
        despliegue sería otra.
  \item \textbf{Automatizar el levantamiento.} Un \textit{pipeline} de
        \textit{scripts} reanudables permitió rehacer corpus y experimentos
        muchas veces sin perder trabajo, y es lo que hace reproducible este
        informe.
\end{itemize}

Qué haría distinto: solicitar la base de eBird en marzo, entrenar la rama de
imagen en la nube desde el inicio para probar arquitecturas mayores, y
organizar al menos una campaña de grabación en terreno para las especies
amenazadas con poco audio chileno.

%-----------------------------------------------------------------------
\section{Trabajo futuro}\label{sec:futuro}

\begin{itemize}[leftmargin=*, itemsep=3pt]
  \item \textbf{Datos para las especies amenazadas.} Solicitar la base
        histórica de eBird (EBD) con esfuerzo y listas completas, que permite
        no detecciones reales en lugar de pseudoausencias. Organizar campañas
        de grabación dirigidas a las trece amenazadas. Explorar
        \textit{embeddings} bioacústicos de modelos entrenados en aves
        (BirdNET, Perch) con clasificadores de pocos ejemplos
        \citep{ghani2023embeddings} para las especies bajo el umbral de
        inclusión.
  \item \textbf{Transformers de visión y rama de imagen en la nube.} Entrenar
        en la nube ConvNeXt o ViT-B con preentrenamiento en ImageNet-21k y
        mayor resolución (el ViT-S/16 a \SI{224}{px} entrenado aquí quedó por
        debajo de la CNN), para comprobar si la imagen sola supera el
        \SI{85}{\percent} y si la atención global ayuda en las especies de
        plumaje casi idéntico. Incorporar TTA en la
        rama visual de la fusión.
  \item \textbf{SDM verdaderamente temporal.} Series de NDVI mensuales de
        Sentinel-2 o MODIS exportadas de forma asíncrona desde GEE, variables
        climáticas del año de observación y validación por bloques espaciales
        y temporales como protocolo por defecto. También cuantificar la
        incertidumbre de los mapas (por ejemplo, con la varianza entre árboles).
  \item \textbf{Monitoreo acústico pasivo.} Evaluar los clasificadores sobre
        paisajes sonoros continuos grabados en terreno, con solapamiento de
        especies y relación señal-ruido baja, que es el escenario de BirdCLEF
        \citep{kahl2023birdclef}, y adaptar el umbral de decisión por especie.
  \item \textbf{Despliegue.} Llevar el prototipo a un servicio estable con
        modelos cuantizados (la CNN de 4 millones de parámetros es candidata a
        correr en dispositivos de campo), registro de las predicciones y
        retroalimentación de expertos para reentrenar.
\end{itemize}

%-----------------------------------------------------------------------
\section{Conclusión general}\label{sec:conclusion}

Este trabajo construyó y evaluó un sistema de monitoreo de biodiversidad
aviar para Chile continental que integra la identificación multimodal de
especies con el modelado de su distribución, a partir exclusivamente de
fuentes públicas de ciencia ciudadana y teledetección. Sus resultados
principales son cuatro. Primero, la fusión de audio e imagen identifica 53
especies con 0,95 de \textit{accuracy}, superando con amplitud la meta del
\SI{85}{\percent}, a cada modalidad por separado y a un modelo de IA generativa multimodal, lo que valida la hipótesis multimodal con datos no
emparejados. Segundo, Random Forest y gradient boosting superan a MaxEnt en
la estimación de distribución con partición aleatoria, pero no al extrapolar a
zonas sin registros, donde todos los modelos empatan. Ni las redes
neuronales ni los parches de cobertura de suelo ni los \textit{embeddings}
de imágenes satelitales aportan, un resultado negativo bien fundamentado que
orienta la práctica. Tercero, el
levantamiento cuantificó la escasez estructural de datos de las especies
amenazadas chilenas y mostró qué estrategias la mitigan en cada modalidad.
Cuarto, todo el proceso quedó en un repositorio reproducible y público
(\url{\repoURL}), con \textit{scripts} reanudables y un prototipo que integra los modelos, que
puede extenderse a otras especies y regiones. Una última lección cruza los dos módulos. El aprendizaje profundo fue
imprescindible donde los datos no tienen estructura tabular: el canto y la
fotografía solo pueden clasificarse aprendiendo representaciones a partir de
la señal cruda, y ahí las redes preentrenadas y su ajuste fino hicieron la
diferencia (31 puntos sobre una regresión logística con rasgos congelados en
audio). Pero cuando los datos ya vienen como una tabla de covariables
ambientales, la estadística clásica sigue siendo muy valiosa. MaxEnt, que en
el fondo es una regresión logística regularizada, obtiene un AUC competitivo
(0,863 frente a 0,884 de Random Forest), empata o gana al extrapolar a zonas
sin registros y entrega algo que los ensambles y las redes no dan con la
misma claridad: coeficientes interpretables que dicen cómo responde cada
especie a la temperatura, la precipitación o la vegetación, y supuestos
explícitos que pueden discutirse con un ecólogo. La elección del método debe
seguir a la naturaleza de los datos y al uso que se dará al modelo, no a la
moda. El sistema no reemplaza al ornitólogo en terreno, pero ofrece una
herramienta concreta para procesar el volumen creciente de registros
ciudadanos y para priorizar dónde y cuándo mirar.
